\oldpage[IV-2]{207}

En el resto de este apartado aplicaremos (7.5.1) al estudio de los
\emph{productos tensoriales completados} $A\widehat\otimes_k B$ de anillos
locales noetherianos que son álgebras sobre un cuerpo $k$.

\begin{lemma}[7.5.5]
\label{IV.2.s7.5.5-fr}
Sean $k$ un cuerpo y $A$, $B$ dos anillos locales noetherianos completos
que contienen $k$, siendo el cuerpo residual de $A$ una extensión finita
de $k$. Sea $C$ el producto tensorial completado
$A\widehat\otimes_k B$ \textup{(\textbf 0$_{\mathrm I}$, 7.7.5)}. Entonces:
\begin{enumerate}
\item[\emph{(i)}] $C$ es un anillo semilocal noetheriano completo.
\item[\emph{(ii)}] $C$ es un $A$-módulo plano y un $B$-módulo plano.
\item[\emph{(iii)}] Si $\mathfrak m$ es el ideal maximal de $A$,
$\mathfrak mC$ está contenido en el radical de $C$, y $C/\mathfrak mC$
es isomorfo a $(A/\mathfrak m)\otimes_k B$.
\item[\emph{(iv)}] Los cuerpos residuales de los componentes locales
de $C$ son extensiones finitas del cuerpo residual de $B$.
\end{enumerate}
\end{lemma}

Las propiedades (i), (iii) y la primera afirmación de (ii) son casos
particulares de (\textbf 0, 19.7.1.2). Para demostrar que $C$ es un
$B$-módulo plano, observemos que, para todo $h>0$,
$C/\mathfrak m^hC=(A/\mathfrak m^h)\otimes_k B$
es un $B$-módulo plano, puesto que $k$ es un cuerpo; basta, por tanto,
aplicar (\textbf 0$_{\mathrm{III}}$, 10.2.6) a cada uno de los componentes
locales de $C$ (de los que $C$ es el producto directo). Por último,
si $\mathfrak n$ es el ideal maximal de $B$, los cuerpos residuales
de los componentes locales de $C$ son también los de los componentes
locales del anillo artiniano
$(A/\mathfrak m)\otimes_k(B/\mathfrak n)$, que, por hipótesis, son
extensiones finitas de $B/\mathfrak n$.
